\documentclass[uplatex,a4paper,twocolumn,10pt]{jsarticle}
\usepackage[margin=20mm]{geometry}
\usepackage{amsmath}
\usepackage{url}
\expandafter\def\expandafter\UrlBreaks\expandafter{\UrlBreaks\do\0\do\1\do\2\do\3\do\4\do\5\do\6\do\7\do\8\do\9}
\title{状態保持型LLMエージェントにおける\newline 情報保持と実行制御の評価}
\author{検討用原稿・著者情報未設定}
\date{}
\begin{document}
\raggedbottom
\maketitle
\noindent\textbf{あらまし案}\quad
履歴要約を既知４課題で比較した. 同時要約は写真で別呼出し要約より総tokenを7.50\%減らし成功したが，CSVは未達だった. 全24実行と974入力の計数から，入力短縮と作成費用，最終成果の関係を調べた.

\noindent\textbf{キーワード}\quad LLMエージェント，ハーネスエンジニアリング，情報保持，実行制御，費用評価

% Working prose started before P120 outcome inspection. Not a submitted manuscript.
\section{問題設定と関連研究}
LLMエージェントが外部APIを用いて複数段階の仕事を行うとき，過去の指示，行動，観測，推論を次の要求へ再入力する費用が蓄積する．必要な情報を残して入力を短縮するには，履歴の保持，状態の引継ぎ，ツール実行を制御するハーネスが必要になる．本研究の目的は，最終成果の性能を維持し，制御と失敗を含めた仕事全体の費用を減らすことである．本稿ではそのうち，一度だけの履歴要約と，その作成費用を取り上げる．

CodeAct\cite{codeact}型のエージェントでは，生成したPythonを実行し，変数を次の工程へ保持できる．一方，変数を保持しても，モデルへ毎回渡す会話履歴の費用は残る．MEM1\cite{mem1}は内部状態の更新と行動生成を結び付け，SKILL.state\cite{skillstate}は構造化状態を更新して履歴の増大を抑える．本稿の通常要約や同時出力を，これらと独立した新方式とは位置付けない．また，局所的なtoken短縮が実際の課金削減に直結しないことや，圧縮後の再取得と行動の変化は既に指摘されている\cite{costproxy}．

本稿は，状態保持型の同じエージェントに全保持，別呼出し要約，同じ応答での要約を実装し，固定予算で全仕事を比較する．要約の要求，実際の採用，次入力の短縮，後続要求全体の入力差，最終成功を区別する．保存した実行列で元履歴を復元する補助計数により，入力短縮と要約生成費用の関係も調べる．対象は既知４課題であり，一般的な品質保証や新しい圧縮原理の提案ではない．

\section{履歴と要約の制御}
\subsection{共通状態と全保持}
元の依頼，公開指示，過去のコード，返却された観測，推論を保存し，Pythonの状態を全工程で維持する．保存用の全履歴と，次のモデル要求へ渡す履歴を分離する．呼出しIDと返却IDの対応，長い観測の表示，予算通知，API実行口は共通とする．全保持では通常の履歴表示を継続する．

\subsection{別呼出し要約}
通常表示の入力が実tokenizerで16,384 token以上，かつ完結したやり取りが３回以上ある最初の境界で，一度だけ同じモデルへ要約を要求する．元のsystem指示と課題，直近２回のやり取りを残し，それより古い完結履歴を要約の対象とする．要約は完了事項，再利用可能な状態，未完了事項，不確実な点の４欄を持つ．各欄は最大32文字列，１文字列は最大2,048文字，compact JSON全体は8,192文字以内とする．構造の検査を通り，通常表示より次入力が厳密に短い場合だけ履歴へ採用する．元履歴とPython状態は削除しない．

\subsection{同じ応答での要約}
別呼出しによる履歴の再入力を避けるため，同じ起動境界で次の\texttt{run\_python}引数に要約欄を付ける．任意欄と，schema及び指示で必須とする欄を別群で比較する．特定toolの強制はせず，\texttt{tool\_choice=auto}を維持する．要約対象は古い完結履歴だけであり，今回のコードの実行結果を先取りして書かせない．今回のコードと実際の返却は，後続履歴に原文で残す．

通常の行動と要約は4,096出力tokenを共有する．有効な基礎引数のコードは元の参照IDと対応付けて一度だけ実行する．要約欄の欠落，不正，打切り，非採用のための再生成はしない．実行後に継続が必要な場合，同じ残予算通知で通常表示と要約表示を数え，短くなった場合だけ採用する．対応を確認した要約metadataは両表示から外すが，原記録には保持する．形式と長さの検査は，内容の真実性や課題への十分性を証明しない．

\section{評価手順と指標}
\subsection{全仕事の比較}
AppWorld\cite{appworld}の写真移動，電気代精算，返金，曲情報のCSV書出しの既知４課題を用いる．全保持対別呼出し，別呼出し対任意欄，別呼出し対必須欄の順に，各群８条件を実行する．各群を事前に固定し，その全結果を閉じてから次群を決めた．各条件は１回であり，同じ４課題の24実行を24個の独立課題とは数えない．未見の予約課題は使わない．

Mac Studio M3 Ultra・256 GiB，AppWorld 0.1.3.post1，MiniMax-M3 IQ4\_XS，llama.cpp b11669を用いる．context上限65,536，１スロット，温度は実行役1.0，要約役0.7，top-p 0.95，環境seed42，モデルseed57から役割別に増加させる．初回はcacheを無効として実測０を確認し，以後は通常のprefix cacheを用いる．課題別の条件順は順方向，逆方向，順方向，逆方向で固定する．

各条件は全役を合わせて80モデル要求，450,000論理token，1,800作業秒までとし，写真160，他80 APIまでとする．要求の上限時間600秒と出力上限4,096は，残予算にも従う．使用量不明の応答は実行せず，API効果と後処理が確認できた場合だけ次の固定条件へ進む．各最終状態を独立評価器で一度判定し，正解や評価結果を実行役へ返さない．失敗，欠測，未実行の枠も残す．

主指標は，全モデル要求についての論理入力と出力の和
\begin{equation}
 T=\sum_j (I_j+O_j)
\end{equation}
である．要約役，追加指示，schema，失敗と回復，cache済み入力も含む．API数，実時間，cache利用量は別に報告し，金額や計算量へ同一視しない．終端使用量の欠測は０や予約上限で補わない．採用判定は全件評価済み，使用量既知，監査一致，課題ごとの成功損失なし，共通成功１件以上，全４課題と共通成功課題の双方で総tokenが少ないことを要求する．必須欄の群では，実際の要約採用１件以上も求める．この開発判定は統計的な非劣性検定ではない．

\subsection{同じ保存行動列での入力差}
全24枠の全モデル要求を対象に，実際の入力と，採用された要約の代わりに元履歴を戻した入力を同じtokenizerで数える．後続のコード，応答，観測，残予算通知，出力上限，schemaを固定する．要約が未採用なら両入力は同じである．別要約の要求は両側に残し，同時出力のmetadataを再入力しない処理も両側へ適用する．

要求$j$での復元入力を$\widetilde I_j$とし，
\begin{equation}
 D=\sum_j(\widetilde I_j-I_j)
\end{equation}
を求める．別呼出し要約の実測入力と出力の和を$C$とし，$D-C$を同じ実行列上の入力差と作成費用の関係として示す．同時出力の要約部分だけの生成tokenは分離測定していないため，独自の$C$を推定しない．終端使用量が不明でも入力だけを計数できる場合があるが，元の総費用の欠測はそのまま残す．

この処理では新たなモデル生成，コード実行，API，課題評価を行わない．元の実測入力tokenとの一致を確認し，仮想入力のcontext上限と仕事予算の超過を明示する．履歴を戻せば実際の行動も変わり得るため，$D$や$D-C$は別方式の総費用，因果効果，品質維持の証拠ではない．直後の短縮量を残り呼出し数へ掛ける近似も行わない．

\section{全仕事の比較結果}
\begin{table*}[t]
\centering\small
\caption{全24条件．$T$は要約役を含む論理token，$M$は全モデル要求，$A$はAPI数．}
\label{tab:whole}
\begin{tabular}{lcrrrcrrr}
\hline
課題 & 基準成否 & 基準$T$ & $M$ & $A$ & 比較成否 & 比較$T$ & $M$ & $A$ \\
\hline
\multicolumn{9}{l}{群1：全保持対別呼出し要約} \\
写真 & 成功 & 135,628 & 16 & 99 & 成功 & 139,871 & 17 & 102 \\
精算 & 成功 & 211,110 & 24 & 25 & 成功 & 211,110 & 24 & 25 \\
返金 & 成功 & 102,087 & 13 & 13 & 成功 & 102,087 & 13 & 13 \\
CSV & 未達 & 358,638$+?_{1}$ & 21 & 53 & 成功 & 440,461 & 32 & 77 \\
\hline
\multicolumn{9}{l}{群2：別呼出し対任意欄の同時要約} \\
写真 & 成功 & 139,871 & 17 & 102 & 成功 & 136,201 & 16 & 99 \\
精算 & 成功 & 211,110 & 24 & 25 & 成功 & 211,110 & 24 & 25 \\
返金 & 成功 & 102,087 & 13 & 13 & 成功 & 102,087 & 13 & 13 \\
CSV & 成功 & 440,461 & 32 & 77 & 未達 & 327,694 & 21 & 80 \\
\hline
\multicolumn{9}{l}{群3：別呼出し対必須欄の同時要約} \\
写真 & 成功 & 139,871 & 17 & 102 & 成功 & 129,374 & 16 & 102 \\
精算 & 成功 & 211,110 & 24 & 25 & 成功 & 211,110 & 24 & 25 \\
返金 & 成功 & 102,087 & 13 & 13 & 成功 & 102,087 & 13 & 13 \\
CSV & 成功 & 440,461 & 32 & 77 & 未達 & 438,073 & 28 & 65 \\
\hline
\end{tabular}
\par\smallskip\footnotesize $+?_{1}$は使用量不明１要求．基準と比較は各群内で別々に実行した．
全保持と任意・必須欄の間の直接比較ではない．
\end{table*}

表\ref{tab:whole}に全24条件を示す．全保持対別呼出しの群では，成功は3/4対4/4であった．写真の別呼出しは次入力を16,425から8,708 tokenへ短縮したが，要約生成に18,455 tokenを使い，全仕事では135,628から139,871へ約3.13\%増えた．実行役だけなら14,212 token減っており，要約役を除くと費用差の符号が逆になる．両側成功の３課題合計も448,825対453,068で約0.95\%増えた．CSVは別呼出しだけ成功したが，全保持の終端１要求の使用量が不明で，全４課題の正確な費用差は未確定である．

任意欄の群は，別呼出し4/4対同時出力3/4成功だった．要約を求めた写真とCSVの２応答とも要約欄を出さず，履歴の要約置換は０件だった．同時出力は777,092 tokenで別呼出し893,529より約13.03\%少ないが，CSVの成功を失った差を含む．入力の圧縮や品質維持の効果と解釈しない．

必須欄の群では２機会とも要約を出し，形式と長さの検査後に採用された．写真は別呼出しと同じく成功し，全仕事は139,871から129,374 tokenへ10,497，約7.50\%減った．APIは両側102，モデル要求は17から16となった．実行役費用は同時出力側で7,958増えたが，別要約18,455の差と合わせて総量が減った．従って専用要約の費用全額が節約されたわけではない．

一方，CSVは同時出力で次入力が21,999から11,110 tokenへ短くなったものの，438,073 tokenで次要求の予算を予約できず未達だった．別呼出しは440,461で成功した．同時出力の作業時間は約1,429.695秒で，別呼出しの約1,304.294秒より長かった．要約欄の欠落や使用量不明による失敗ではないが，終端理由だけから要約の意味的欠落が原因とは確定できない．

必須欄の全４課題は880,644対893,529 token，両側成功３課題は442,571対453,068であった．成功損失１を含む全体を性能維持の削減とはしない．精算と返金は全群で要約が起動せず，方式間のtoken差も０だった．この２件を圧縮の成功例へ数えない．固定した採用基準は３群とも未達であった．

各群の最初の介入前の完全要求は写真14位置，CSV12位置で全て一致した．介入のない精算24位置と返金13位置も一致した．この照合は，後続の行動や費用の不変，個別の要因の因果効果までは示さない．全８枠とserverを閉鎖してから成否と費用を確認し，結果を見た条件変更や失敗分の補充実行は行っていない．

\section{入力短縮と作成費用の関係}
全487保存要求の２表示，974入力を計数し，欠測と不一致は０だった．使用量既知486要求では元の実測入力に一致した．不明１要求の入力も数えられたが，元の終端使用量は不明のまま保持する．計数本体は37.420秒，モデル生成，API，評価の追加は０だった．

\begin{table}[t]
\centering\small
\caption{要約採用後の保存入力差．$D-C$は実測総節約ではない．}
\label{tab:inputs}
\begin{tabular}{llrrr}
\hline
方式 & 課題 & $D$ & 別要約$C$ & $D-C$ \\
\hline
別要約 & 写真 & 15,434 & 18,455 & $-3,021$ \\
別要約 & CSV & 200,507 & 21,419 & 179,088 \\
必須欄 & 写真 & 8,023 & -- & -- \\
必須欄 & CSV & 163,335 & -- & -- \\
\hline
\end{tabular}
\par\smallskip\footnotesize 別要約の値は３群でそれぞれ同じ．必須欄は群3．
同時出力の要約だけの出力費用は未分離であり，--を０とはしない．
\end{table}

表\ref{tab:inputs}の写真の別要約では，要約後に２要求があり，各入力差7,717の合計は作成費用に3,021足りなかった．CSVの別要約では３要求目で累積入力差が31,659となり，作成費用21,419を上回った．終了までの19要求の差は200,507だった．全群の精算と返金及び群2の任意欄は要約未採用で，入力差０だった．

必須欄の写真では要約後の１要求で8,023，CSVでは15要求で163,335の入力差があった．後者で入力短縮が実際に起きても，元の未達は変わらない．CSV別要約の実測440,461に$D$を足した640,968，必須欄438,073に$D$を足した601,408は仕事予算450,000を超える．これらを実行可能な全保持の比較結果としない．復元入力と実際の出力予約がcontext上限を超える例はなかった．

\section{考察と適用範囲}
写真の必須欄条件では，通常の別呼出し要約に対して成功と総token削減が両立した．しかし，実際に要約を使ったもう１件のCSVでは，入力が短くなっても成功を失った．要約の生成，入力への採用，総費用，最終成果を別々に測る必要がある．要約の作成費用だけを除いた比較や，成功した課題だけの合計は，全条件の性能維持を示さない．

対象は開発で使用した既知４課題，１モデル，各条件１回である．写真の7.50\%はこの実行対の観測値であり，未見課題の平均削減率や統計的な品質の非劣性ではない．３群は前の結果を見て設計した順次開発である．要約の作成時点，温度，指示，schema，出力配分と後続行動の差を含む方式全体の比較であり，同時生成だけの効果を分離していない．

全24枠の入力再構成も，実際の行動列と通知を固定した事後診断である．仮想入力を使った再生成と評価はしておらず，失われた意味の十分性を検証していない．同じseedでもcacheや数値実行，ホスト負荷の影響が残り，論理tokenを課金，FLOPs，実時間と同一視しない．完全履歴，課題データ，私有推論を公開せず，全条件の集約，固定条件，ソースと監査記録を公開する\footnote{\url{https://github.com/sophiria1225/AI-}}．

\section{結論}
既知４課題の履歴制御を３群，計24実行で比較した．必須欄の同時要約は写真で別呼出し要約より総tokenを7.50\%減らして成功したが，CSVの成功を失った．974入力の補助計数では，写真の別要約は後続２要求で作成費用を回収せず，CSVでは入力差が大きくても同時要約の成功を保証しなかった．全条件での品質維持と削減を確立した結果ではない．今後は，必要な状態を残す規則と作成時点を独立した課題で検証する必要がある．

\begin{thebibliography}{99}
\raggedright
\bibitem{codeact} X. Wang, Y. Chen, L. Yuan, Y. Zhang, Y. Li, H. Peng, and H. Ji, ``Executable code actions elicit better LLM agents,'' Proc. ICML, PMLR, vol.235, pp.50208--50232, Vienna, Austria, July 2024.
\bibitem{appworld} H. Trivedi, T. Khot, M. Hartmann, R. Manku, V. Dong, E. Li, S. Gupta, A. Sabharwal, and N. Balasubramanian, ``AppWorld: a controllable world of apps and people for benchmarking interactive coding agents,'' Proc. ACL, vol.1, pp.16022--16076, Bangkok, Thailand, Aug. 2024.
\bibitem{mem1} Z. Zhou, A. Qu, Z. Wu, S. Kim, A. Prakash, D. Rus, J. Zhao, B. K. H. Low, and P. P. Liang, ``MEM1: Learning to synergize memory and reasoning for efficient long-horizon agents,'' arXiv:2506.15841v1, June 2025.
\bibitem{skillstate} S. Badhe, P. Tiwari, and J. Chung, ``SKILL.state: Scalable long-horizon agent skills,'' arXiv:2608.26263v1, Aug. 2026.
\bibitem{costproxy} S. Weinberger and A. Hozez, ``Token reduction is not cost reduction,'' arXiv:2607.12161v5, Aug. 2026.
\end{thebibliography}


\clearpage
\onecolumn
\noindent\textbf{Title}\quad Evaluating Information Retention and Execution Control in Stateful LLM Agents

\noindent\textbf{Authors}\quad 未設定

\noindent\textbf{Keywords}\quad LLM agents, harness engineering, information retention, execution control, cost evaluation

\noindent\textbf{Summary draft}\quad We compare checkpoint policies across four known tasks. Required inline summaries reduced total tokens by 7.50\% on the photo task relative to separate summaries, while preserving success, but failed on CSV export. An analysis of 24 runs and 974 inputs separates input reduction, summary cost, and task outcomes.
\end{document}
